% =====================================================================
%  Chapter 4 — Results
%  Section 4.1 written from the exploratory analysis.
%  Style notes: key terms in \textbf so the chapter can be skimmed;
%  enumerations use itemize rather than running prose.
% =====================================================================

\section[Dataset Characterisation and Exploratory Analysis]{Dataset Characterisation and Exploratory\\ Analysis}
\label{sec:res_eda}

This section reports what the dataset contains before any model is trained. It is
presented as a result rather than as preparatory work, because its findings
\textbf{constrained the design of everything that follows}: they
changed which question could honestly be asked, how any answer must be
validated, and which features can be trusted.

\subsection{Cohort composition}

Figure~\ref{fig:balance} shows the composition. The class imbalance is mild,
\textbf{33 control against 25 PD}, or roughly
$57\,\%$ to $43\,\%$. It is small enough not to require resampling, but large
enough that plain accuracy would mislead: a classifier that always predicted the
majority class would score $57\,\%$ while being diagnostically useless. This
motivates the use of \textbf{area under the ROC curve} as the headline metric,
supported by balanced accuracy, sensitivity and specificity.

\begin{figure}[H]
  \centering
  \includegraphics[width=0.92\textwidth]{eda_class_balance.png}
  \caption[Cohort composition]{Subjects per group, and the number of recordings
  contributed by each subject.}
  \label{fig:balance}
\end{figure}

% ---------------------------------------------------------------------
\subsection{Recording duration differs between groups, and the chair is why}
\label{sec:res_duration}

PD subjects take longer to complete the test, and the difference is
\textbf{concentrated almost entirely in the variant that includes the chair},
as Figure~\ref{fig:duration} and Table~\ref{tab:duration} show.

\begin{figure}[H]
  \centering
  \includegraphics[width=0.95\textwidth]{eda_duration_confound.png}
  \caption[Recording duration by group and variant]{Distribution of recording
  durations for each group, split by test variant. Dashed lines mark each
  group's median.}
  \label{fig:duration}
\end{figure}

\begin{table}[H]
  \centering
  \caption[Recording duration by group and variant]{Recording duration by group and test variant. $\delta$ is Cliff's
  delta; $p$ is the Mann-Whitney probability value.}
  \label{tab:duration}
  \begin{tabular}{@{}lccccc@{}}
    \toprule
    \textbf{Variant} & \textbf{Control median} & \textbf{PD median} &
    \textbf{Difference} & $\boldsymbol{p}$ & $\boldsymbol{\delta}$ \\
    \midrule
    \textbf{test1} (with chair)    & 9.17~s & 10.09~s & $+10.1\,\%$ & $0.0016$ & $0.278$ \\
    \textbf{test2} (without chair) & 8.04~s & 8.31~s  & $+3.4\,\%$  & $0.141$  & $0.131$ \\
    \bottomrule
  \end{tabular}
\end{table}

In test1 the gap is unlikely to be coincidence and of small-to-moderate
magnitude. In test2 it falls to a negligible effect size entirely consistent
with ordinary between-subject variation. \textbf{Removing the chair removes most
of the difference}, which is the behaviour predicted by the clinical picture:
rising from a chair is among the movements most affected by bradykinesia
\cite{postuma2015}, so the additional seconds accumulate in the sit-to-stand and
stand-to-sit transitions rather than during steady walking. The same comparison
was reproduced through an independently written analysis path and reached the
same conclusion.

Two consequences follow, and both are carried through the rest of this thesis:

\begin{itemize}
  \item \textbf{Elapsed time is a confound, not a feature.} Any method that
        responds to recording length will appear to detect the disease. A
        \texttt{duration\_only} control model is therefore evaluated alongside
        every other feature set, so that its score acts as a floor below which no
        result can be considered evidence of gait.
  \item \textbf{The chair-free variant, test2, is the cleaner comparison.} With the chair removed the two
        groups take almost the same time, so a difference found there cannot be
        explained by one group simply spending longer in front of the radar.
\end{itemize}

% ---------------------------------------------------------------------
\subsection{Nearly every feature differs, which is a warning rather than a result}
\label{sec:res_effects}

Comparing the two groups across all 25 features gives a result that appears
excellent: \textbf{24 of the 25 show a difference unlikely to be coincidence}
after correction for multiple comparisons. It is in fact a warning sign.

It is not plausible that a walk differs between the groups in 24 independent
ways. It is very plausible that \textbf{one shared factor differs} and that it
surfaces in nearly every measurement, and Section~\ref{sec:res_duration} has
already identified that factor. Figure~\ref{fig:effects} confirms the diagnosis
directly.

\begin{figure}[H]
  \centering
  \includegraphics[width=0.95\textwidth]{eda_effect_sizes.png}
  \caption[Effect size per feature]{Cliff's delta for each of the 25 features,
  sorted. Red marks the features identified in advance as growing with recording
  length; blue marks the remainder. Dotted guides mark the negligible, small and
  medium thresholds.}
  \label{fig:effects}
\end{figure}

Two observations stand out:

\begin{itemize}
  \item \textbf{The five largest effects are all length-dependent.} Ranked by
        magnitude, the top five are all summed-energy features, and the first
        length-independent feature ranks sixth at $\delta = 0.41$. The strongest
        apparent differences therefore come from features that partly encode
        elapsed time rather than the shape of the motion.
  \item \textbf{No feature reaches a large effect at all.} The largest value in
        the dataset is $\delta = 0.461$, marginally below the conventional
        threshold for a large effect~\cite{romano2006}, and it is one of the
        length-dependent features.
\end{itemize}

\noindent The general lesson is that \textbf{a large number of significant
results is weak evidence, not strong evidence}: it indicates that the
measurements are not independent. Two or three differing features would have
been more convincing, because that pattern would suggest something specific
rather than something that touches every measurement.

\begin{figure}[H]
  \centering
  \includegraphics[width=0.95\textwidth]{eda_top_feature_distributions.png}
  \caption[Group distributions of the strongest length-independent features]{
  Distribution of each group's values for the strongest length-independent
  features. Wider regions indicate more recordings at that value.}
  \label{fig:distributions}
\end{figure}

Figure~\ref{fig:distributions} shows what a probability value cannot. A small
probability establishes that the groups differ \textbf{on average}; it says
nothing about whether an \textbf{individual} could be assigned to the correct
group, which is the only question a diagnostic tool has to answer. The
distributions answer it directly: even for the best features, the two groups
\textbf{overlap almost completely}. The value read from a PD subject is a value
that many control subjects also produce.

This observation sets the expectation for the modelling results reported later
in this chapter. Those results are modest, and \textbf{this overlap is the
reason}. It was visible in the dataset before any model was built, which makes it
a property of the problem rather than a shortcoming of a particular method.

% ---------------------------------------------------------------------
\subsection{A walk identifies the walker}
\label{sec:res_fingerprint}

Every subject walks the same 3~m corridor, so the following analysis concerns
similarity rather than physical distance. Using the dissimilarity of
Equation~\eqref{eq:dist}, each subject yields two numbers: the
\textbf{within-subject} distance among their own recordings, and the
\textbf{between-subject} distance to other subjects' recordings. If a walk
carried no personal signature, the two would be equal.

\begin{figure}[H]
  \centering
  \includegraphics[width=0.95\textwidth]{11_within_vs_between.png}
  \caption[Subject fingerprint strength]{One point per subject. The horizontal
  position is the mean distance among that subject's own recordings; the
  vertical position is the mean distance to ten other subjects' recordings drawn
  at random. The dashed line marks equality. Four subjects are ringed and named
  to indicate the range observed.}
  \label{fig:fingerprint}
\end{figure}

The result, shown in Figure~\ref{fig:fingerprint}, is unambiguous. Every one
of the \textbf{58 subjects lies above the
line}, with no exceptions. A subject's own recordings sit a median distance of
$1.81$ apart, against $4.59$ to other subjects' recordings, so a subject's own
walks are approximately \textbf{2.5 times more alike} than anyone else's.

The interpretation requires care. The quantity being tested is not consistency:
the subject with the widest spread among their own recordings, at $6.87$, still
lies above the line because they sit $9.34$ from everyone else. What makes a
subject identifiable is being \textbf{more similar to themselves than to
others}, which is a comparison rather than a property of an individual.

This is a \textbf{hazard rather than a useful finding}. Were one recording of a
PD subject used for training and a second recording of the same subject used for
evaluation, a model would not need to learn anything about Parkinson's disease:
it could recognise \textbf{the individual} and return the label already
associated with them. The resulting score would be genuine and the conclusion
drawn from it would be false. Two rules follow, and both are applied throughout:

\begin{itemize}
  \item Every subject is held out \textbf{in full}, never one recording at a
        time, which is the \textbf{leave-one-subject-out} protocol of
        Section~\ref{sec:validation};
  \item Uncertainty is estimated by resampling \textbf{whole subjects} rather
        than individual recordings, because six recordings of one subject are
        not six independent observations and treating them as such would
        understate the uncertainty.
\end{itemize}

% ---------------------------------------------------------------------
\subsection{How the exploratory analysis constrained the design}
\label{sec:res_eda_impact}

Three elements of the original plan had to be revised once the dataset had been
examined. Table~\ref{tab:eda_impact} records them, because the revisions are
themselves a contribution of this thesis.

\begin{table}[!htbp]
  \centering
  \caption[Design decisions revised by the exploratory analysis]{Design decisions revised by the exploratory analysis.}
  \label{tab:eda_impact}
  \begin{tabular}{@{}L{4.2cm}L{4.2cm}L{6.2cm}@{}}
    \toprule
    \textbf{Planned} & \textbf{Found} & \textbf{Revised} \\
    \midrule
    Measure how accurately a model separates Parkinson's from healthy control. &
    PD subjects take longer to complete the test. &
    A model responding only to elapsed time would score well having learned
    nothing about gait. The objective became measuring \textbf{how much
    separation survives once duration is removed}. \\
    Split the recordings at random into training and test sets. &
    Two recordings of the same subject are far more alike than recordings of two
    different subjects. &
    A model could score well by \textbf{recognising the individual}. Splitting is
    performed \textbf{by subject}, never by recording. \\
    Decide by inspection which features grow purely with recording length, and
    exclude them. &
    Those judgements were wrong in both directions; several features assumed to
    be safe track duration closely. &
    Membership of the length-independent set is \textbf{determined by testing
    each feature against measured duration}. \\
    \bottomrule
  \end{tabular}
\end{table}

The third revision has a cause worth stating, because it is not obvious. A
longer recording is not the same walk performed more slowly. It contains
proportionally \textbf{more standing, hesitating and turning} and proportionally
less steady walking, so the mixture of activity inside it differs. A mean, which
does not accumulate with length in the way a sum does, therefore still shifts
when a recording is longer, which is exactly why membership of the
length-independent set cannot be settled by inspection.

% =====================================================================
\section{The Null Floors}
\label{sec:res-floors}

Every model in this chapter is judged against the two deliberately
impoverished predictors of Table~\ref{tab:floors}, scored under the identical
leave-one-subject-out protocol of Section~\ref{sec:validation}. Any model not
clearly above them has demonstrated nothing about gait.

\begin{table}[htbp]
  \centering
  \caption[The null floors]{The null floors.}
  \label{tab:floors}
  \small
  \begin{tabular}{@{}lcp{8.2cm}@{}}
    \toprule
    Floor & AUC & How it is computed \\
    \midrule
    Recording duration alone & 0.610 &
      One number per subject, how long their recordings last, fed to a
      logistic regression under the same leave-one-subject-out loop; the
      \texttt{duration\_only} cell of Table~\ref{tab:aucgrid}. \\
    Window count alone & 0.620 &
      No model at all: subjects are ranked directly by how many windows they
      yield (a restatement of duration, Section~\ref{sec:preproc}) and the
      AUC of that ranking is computed. \\
    \bottomrule
  \end{tabular}
\end{table}

Both floors carry \textbf{no gait content whatsoever}, yet both sit above
0.6. They quantify, before any model is discussed, how much apparent
performance the recording-length confound of Section~\ref{sec:res_duration} supplies for free.

% =====================================================================
\section{Classical Baselines}
\label{sec:res-baselines}

Table~\ref{tab:aucgrid} reports the full grid promised in
Section~\ref{sec:baselines}: three classifiers by five feature sets, each
cell the subject-level AUC of a complete leave-one-subject-out run, drawn as
a heat map in Figure~\ref{fig:aucgrid}.

\begin{table}[htbp]
  \centering
  \caption[Classical baseline AUCs]{Classical baseline AUCs: three classifiers by five feature sets
  (feature sets as defined in Table~\ref{tab:featuresets}).}
  \label{tab:aucgrid}
  \small
  \begin{tabular}{@{}lccccc@{}}
    \toprule
    Model & \texttt{\scriptsize all\_features} &
      \texttt{\scriptsize duration\_invariant} &
      \texttt{\scriptsize shape\_clean} &
      \texttt{\scriptsize confounded\_only} &
      \texttt{\scriptsize duration\_only} \\
    \midrule
    Logistic regression & 0.633 & 0.627 & 0.621 & \textbf{0.675} & 0.610 \\
    SVM                 & 0.623 & 0.601 & 0.577 & 0.634 & 0.287\textsuperscript{*} \\
    Random forest       & 0.605 & 0.601 & 0.583 & 0.618 & 0.450 \\
    \bottomrule
  \end{tabular}

  \vspace{4pt}
  \parbox{0.92\textwidth}{\footnotesize \textsuperscript{*}A one-feature
  kernel model can score below chance under leave-one-subject-out because
  its probability calibration is refitted in every fold, which makes the
  pooled scores incomparable across folds and scrambles the subject ranking;
  the anomaly is a calibration artifact of that cell, not a property of the
  data, as the logistic regression on the same single feature (0.610)
  shows.}
\end{table}

\begin{figure}[H]
  \centering
  \includegraphics[width=0.9\textwidth]{baseline_auc_grid.png}
  \caption[The baseline grid as a heat map]{The grid of Table~\ref{tab:aucgrid} as a heat map. Every row peaks
  in the \texttt{confounded\_only} column: the warmth gradient runs toward
  the features that encode recording length, not away from them.}
  \label{fig:aucgrid}
\end{figure}

The logistic regression row is read \textbf{right to left}. Its
\texttt{duration\_only} cell, recording length alone, scores 0.610. Its
\texttt{confounded\_only} cell, the measurements that merely encode length,
scores \textbf{0.675}, the best cell in the table. Across the four
multi-feature sets the ordering follows the duration content exactly:
0.675, then 0.633, 0.627 and 0.621 as duration information is progressively
removed. \textbf{The more duration information the features carry, the
better the classifier does}: the grid is the controlled experiment
Section~\ref{sec:baselines} designed, and it came out in the direction that
indicts the confound.

The bootstrap confidence intervals of Section~\ref{sec:confounds} say which
of these numbers is real. The interval on the confounded cell runs from
\textbf{0.527 to 0.813} and excludes 0.5: the length signal is genuinely
there. The interval on the clean cell runs from \textbf{0.473 to 0.770} and
contains 0.5: a gait signal has \textbf{not been shown to exist}.

% =====================================================================
\section{Deep Models}
\label{sec:res-deep}

Table~\ref{tab:deepauc} reports the three models of
Section~\ref{sec:deep} under the full nested protocol, and
Figure~\ref{fig:campaign} places them, together with two of the exploratory
variants of Section~\ref{sec:res-ablations} (foot channel only and turn
excluded), against the floors.

\begin{table}[H]
  \centering
  \caption[Deep models under leave-one-subject-out validation]{Deep models under leave-one-subject-out validation.}
  \label{tab:deepauc}
  \small
  \begin{tabular}{@{}lcc@{}}
    \toprule
    Model & Pooled AUC & 95\,\% CI \\
    \midrule
    SmallCNN        & 0.616 & [0.468, 0.760] \\
    ResNet-18 probe & 0.654 & [0.509, 0.798] \\
    Envelope-LSTM   & 0.560 & [0.404, 0.714] \\
    \bottomrule
  \end{tabular}
\end{table}

\begin{figure}[H]
  \centering
  \includegraphics[width=\textwidth]{res_campaign.png}
  \caption[The deep campaign against the floors]{The three models and two of the exploratory input variants. Dots are
  pooled subject-level AUCs, bars their 95\,\% confidence intervals; the
  vertical lines are chance and the two floors of Table~\ref{tab:floors}.}
  \label{fig:campaign}
\end{figure}

Three observations carry the section:

\begin{itemize}
  \item \textbf{Every confidence interval contains the duration floor.} No
        configuration is shown to exceed a predictor that reads only how long
        the walk took. Two intervals (the probe and the foot-only variant)
        clear chance; none clears the floor.
  \item \textbf{More capacity and transfer learning do not change the
        picture.} The scratch-trained SmallCNN (23{,}682 parameters), the
        ImageNet probe (1{,}026 trainable), and the sequence model see the
        problem through entirely different representations and land in the
        same band.
  \item \textbf{The representation that cannot see image texture scores
        lowest.} The Envelope-LSTM, which receives velocity curves rather
        than images and is therefore blind to per-recording texture, reaches
        only 0.560, and in 28 of its 58 folds no epoch beyond the first
        improved on the first epoch's result. The most economical reading is
        that what the image models add over it is texture rather than gait,
        a reading the attention analysis of Chapter~\ref{ch:discussion}
        tests directly.
\end{itemize}

% =====================================================================
\section{Ablation Studies}
\label{sec:res-ablations}

Table~\ref{tab:ablations} reports four input variants, each a re-run of the
SmallCNN under the full pipeline with exactly one component changed
(Section~\ref{sec:preproc}; the raw-magnitude variant instead bypasses the
enhanced representation of Section~\ref{sec:dataset}). Each answers a design
question raised in the methods chapter, and the first two also appear in
Figure~\ref{fig:campaign}.

\begin{table}[htbp]
  \centering
  \caption[Exploratory input variants]{Exploratory input variants (full protocol, pooled AUC).}
  \label{tab:ablations}
  \small
  \begin{tabular}{@{}lcp{7.2cm}@{}}
    \toprule
    Variant & AUC & What it answers \\
    \midrule
    Foot channel only & 0.653 &
      The source publication's prediction that the foot channel carries more
      usable signal. A first pass reached 0.674; the pre-registered re-test
      settled at 0.653, indistinguishable from the two-channel models. No
      channel advantage survived. \\
    Turn excluded ($n=54$) & 0.568 &
      Whether the signal lives in steady walking. Restricting to
      single-direction windows lowers the score toward chance: the models
      were reading the turn-and-stand content, where the group time
      difference of Section~\ref{sec:res_duration} lives. \\
    Per-window standardisation & 0.647 &
      Whether absolute brightness hides a signal or a nuisance. Deleting it
      removes most of the length residue and uncovers nothing beneath:
      nuisance. \\
    Raw $|$STFT$|$ & 0.661 &
      What the contrast enhancement contributes. Bypassing it restores
      absolute energy and the scores track recording length twice as
      strongly: the rise is the nuisance channel, not gait. \\
    \bottomrule
  \end{tabular}
\end{table}

In total, eleven deep configurations were explored across the campaign and
these variants. Those that inform a conclusion are reported in this chapter,
and the complete run record is archived with the project, so no configuration
was selected for its score.

% =====================================================================
\section{Decision Performance at a Fixed Threshold}
\label{sec:res-decisions}

AUC measures ranking; a screening deployment must decide. Following
Section~\ref{sec:metrics}, each model's per-subject probabilities are cut at
the prevalence of the positive class ($25/58 = 0.431$), and
Table~\ref{tab:decisions} reports the resulting decisions for the classical
baseline and the three deep models, and Figure~\ref{fig:confusion} draws the
four confusion matrices.

\begin{table}[htbp]
  \centering
  \caption[Decision metrics at the prevalence threshold]{Decision metrics at the prevalence threshold (positive class:
  PD).}
  \label{tab:decisions}
  \small
  \begin{tabular}{@{}lccccc@{}}
    \toprule
    Model & Recall & Specificity & Precision & F1 & Bal.\ acc. \\
    \midrule
    Logistic regression (\texttt{shape\_clean}) & 0.68 & 0.55 & 0.53 & 0.60 & 0.61 \\
    SmallCNN                        & 0.60 & 0.67 & 0.58 & 0.59 & 0.63 \\
    ResNet-18 probe                 & 0.68 & 0.48 & 0.50 & 0.58 & 0.58 \\
    Envelope-LSTM                   & 0.60 & 0.64 & 0.56 & 0.58 & 0.62 \\
    \bottomrule
  \end{tabular}
\end{table}

\begin{figure}[htbp]
  \centering
  \includegraphics[width=\textwidth]{res_confusion.png}
  \caption[Confusion matrices at the threshold]{Subject-level confusion matrices at the prevalence threshold, one
  per model: counts and row-normalised percentages.}
  \label{fig:confusion}
\end{figure}

\begin{figure}[htbp]
  \centering
  \includegraphics[width=0.95\textwidth]{res_subject_scatter.png}
  \caption[Per-subject probabilities, sorted]{Per-subject probabilities of the ResNet-18 probe, sorted. The
  violet line is the decision threshold at 0.431, the cohort's own PD share;
  an X marks each misclassified subject.}
  \label{fig:scatter}
\end{figure}

At this threshold the best deep model catches \textbf{17 of the 25 PD
subjects}, misses 8, and falsely flags \textbf{17 of the 33 controls};
balanced accuracy spans 0.58 to 0.63 across the four models. The per-subject
view of Figure~\ref{fig:scatter} shows why no other threshold would do
better: the two classes are interleaved along the entire probability ramp,
so sliding the line only trades one kind of error for the other. The weak
decisions are a property of the scores, not of the cut.

\begin{figure}[htbp]
  \centering
  \includegraphics[width=\textwidth]{res_case_gallery.png}
  \caption[One case per confusion cell]{One subject per confusion cell for the two image models: the most
  confident correct PD, the most confidently missed PD, the most confident
  false alarm, and the most confident correct control, each shown through a
  representative window (foot channel) with the model's probability.}
  \label{fig:casegallery}
\end{figure}

Figure~\ref{fig:casegallery} attaches a concrete case to each cell of the
confusion matrix, and the pattern is legible in the backgrounds rather than the
strides: for both models, the subjects \textbf{called PD} (correctly or not)
sit on bright, noisy backgrounds, and the subjects \textbf{called control}
(correctly or not) on dark, clean ones. The most extreme case is the
probe's false alarm, a control flagged at 0.81 whose window is almost pure
noise texture, beside a genuine PD subject waved through at 0.24 on a clean
background. The errors are not random; they are the systematic price of a
score that follows recording character rather than gait, an observation
taken up in the discussion.
